\newcommand{\usource}{u_\text{src}}
\newcommand{\uslm}{u_\text{slm}}
\newcommand{\uz}{u_\text{z}}
\newcommand{\utarget}{u_\text{target}}
\newcommand{\upupil}[1]{u_{\text{pupil#1}, p}}
\newcommand{\Fuslm}{\mathcal{U}_\text{slm}}
\newcommand{\Futarget}{\mathcal{U}_\text{target}}
\newcommand{\Fupupil}[1]{\mathcal{U}_{\text{pupil#1}, p}}
\newcommand{\filter}{\mathcal{P}}
\newcommand{\pupil}{\mathcal{M}}
\newcommand{\fourier}{\mathcal{F}}
\newcommand{\transfer}{\mathcal{H}}
\newcommand{\invfourier}{\mathcal{F}^{-1}}
\newcommand{\fx}{f_x}
\newcommand{\fy}{f_y}
\newcommand{\cx}[1]{c_{x, #1}}
\newcommand{\cy}[1]{c_{y, #1}}
\newcommand{\numpixelxslm}{N_{x, slm}}
\newcommand{\numpixelyslm}{N_{y, slm}}
\newcommand{\numpixelxdmd}{N_{x, dmd}}
\newcommand{\numpixelydmd}{N_{y, dmd}}

In this section, we first review the conventional single-source holographic image formation model before introducing the multi-source image formation model of our system.

\begin{figure}[t]
    \centering
    \includegraphics[width=1.0\linewidth]{./figures/schematic_conceptual.pdf}
    \caption{\textbf{System Architecture.} Multiple mutually incoherent sources illuminate a fast phase-only SLM, creating a high-\'etendue backlight. An additional amplitude display is placed at the Fourier plane to remove ghost image artifacts created by HDOs and the multisource illumination. The phase and amplitude patterns are optimized jointly for a target light field in a content adaptive manner. We illustrate our design following the style of the schematic in \cite{grace2023multi} for easier comparison.}
\label{fig:schematic_conceptual}
\end{figure}

\subsection{Single-Source Holographic Image Formation Model} 

 % with time-multiplexing introduced by \cite{choi2022flex}. 
For on-axis Fresnel holography, a collimated beam from a laser source illuminates an SLM with a normally incident, coherent field $\usource$. The SLM imparts a spatially varying phase delay $\phi$ to the field which propagates a distance $z$ along the optical axis. The wavefront at this plane can be mathematically described using the angular spectrum method (ASM)~\cite{goodman2005fourier} as a function of the phase pattern and distance from the SLM:

\begin{equation}
\begin{split}
    & f \left(\phi, z \right) = \invfourier \left\{ \fourier \left\{ e^{i\phi \left(x, y \right)} \usource \left(x, y \right) \right\} \cdot \transfer \left(\fx, \fy; z \right) \right\} \\
    &\transfer \left(\fx, \fy; z \right) = \begin{cases}
    e^{i\frac{2\pi}{\lambda} z \sqrt{1- \left( \lambda \fx \right) ^2- \left( \lambda \fy \right) ^2} } & \text{if} \, \sqrt{\fx^2+\fy^2} < \frac{1}{\lambda}.\\
    0 & \text{otherwise}.
\end{cases}
\end{split}
\label{eq:asm}
\end{equation}
%
Here, $\lambda$ is the wavelength of light, $x, y$ are the spatial coordinates on the SLM, $\fx, \fy$ are the frequency coordinates, %$\Fuslm$ is the fFourier transform of the single-source modulated field at the SLM, 
and $\transfer$ is the transfer function of the ASM. The operator $f$ models free-space propagation between the parallel SLM and target planes separated by a distance $z$. For notational convenience, we omit the dependence
of the fields on $x$ and $y$. The intensity generated by a holographic display at a distance $z$ in front of the SLM is therefore $\left| f \left(\phi, z \right) \right|^2$. If a high-speed SLM is available, a time-multiplexed variant of the image formation is $\sum_{t=1}^T \left| f \left(\phi^{(t)}, z \right) \right|^2 \!\! / T$, where $T$ phase SLM patterns $\phi^{(t)}, t = 1, \ldots, T$ are rapidly displayed in sequence, and the resulting intensities are averaged by the users' eye~\cite{choi2022flex}.


\subsection{Multi-Source Holographic Image Formation Model with Fourier Modulation}
To extend the single-source image formation model to our system, we modify the formulation to incorporate off-axis collimated illumination traveling in direction
$\mathbf{k} = (k_x, k_y, k_z)$ and a programmable amplitude mask $\filter$ at the Fourier plane of the holographic display system. This results in the model

\begin{align}
    & f^{(j)} \left(\phi, \filter, z \right) = 
    \invfourier \left\{ \Fuslm^{(j)} \left(\fx, \fy; \phi \right) \cdot \transfer \left(\fx, \fy; z \right) \cdot \filter(\fx, \fy) \right\} \nonumber \\     
    & \Fuslm^{(j)} \left(\fx, \fy ; \phi \right) = \fourier \left\{ e^{i\phi \left(x, y \right)} \usource^{(j)} \left(x, y \right) e^{i \, \mathbf{k}^{(j)} \cdot \mathbf{x}}\right\}
\label{eq:multi}
\end{align}
%
where $j$ is the index of the source, 
%$\theta_i$ is the propagation direction of the collimated light source $i$, 
and $\usource^{(j)} \left(x, y \right)$ is the complex-valued field modeling any deviations in amplitude and phase of source $j=1, \dots, J$ from a perfect plane wave $e^{i \, \mathbf{k}^{(j)} \cdot \mathbf{x}}, \mathbf{x} = (x,y,z)$. Moreover, $\usource^{(j)} \left(x, y \right)$ can optionally also include per-source, time-dependent modulation, such as switching individual lasers on and off. In our setup, we do not consider this case and assume that all sources are turned on at all times. Please refer to the supplemental material for more discussions about the generalized configuration with amplitude-controllable laser sources $\usource^{(j)} \left(x, y \right)$.



\subsection{Stochastic Optimization of Light Field Holograms}

To reconstruct a light field, we use gradient descent to optimize a set of time-multiplexed phase patterns $\phi^{(t)}$ and corresponding Fourier masks $\filter^{(t)}$ by minimizing the following objective:
\begin{equation}
    \underset{\left\{\phi^{(t)}, \filter^{(t)} \right\}}{\textrm{minimize}} \left  \lVert s \sqrt{\frac{1}{T}\sum_{t=1}^{T} \sum_{j=1}^{J} |\text{H2LF}\left( f^{(j)} \left(\phi^{(t)}, \filter^{(t)}, z\right) \right)}|^2 - \textrm{lf}_\text{target} \right \rVert 
\label{eq:h2lf}
\end{equation}
where $s$ is a scale factor, $\textrm{lf}_\text{target}$ is the amplitude of the target light field, and $\text{H2LF}$ is a hologram-to-light field transformation, such as the Short-Time Fourier Transform (STFT)~\cite{zhang2009wigner,padmanaban2019olas}. 

The memory consumption of the above optimization problem is huge due to time multiplexing, multiple sources, and the explicit modeling of HDOs. Therefore, it is impractical to realize $\text{H2LF}$ using the STFT since it reconstructs a whole light field and the memory consumption would explode for a dense $\text{lf}_\text{target}$. To solve this problem, we devise a $\textit{stochastic}$ version of Eq. \ref{eq:h2lf} that allows us to optimize a single light-field view rather than a full light field in each iteration of the optimization routine. 

For this purpose, we randomly chose a view $p$ of the target light $\text{lf}^{(p)}_\text{target}$ in each iteration and run a gradient descent step of Eq.~\ref{eq:h2lf}. A binary pupil mask $\pupil^{(p)}$ in the Fourier plane in the hologram-to-light field transform is applied to reconstruct one specific view as 
%
\begin{equation}
\begin{split}
    & \text{H2LF}^{(p)}\left( f^{(j)} \left(\phi^{(t)}, \filter^{(t)}, z\right) \right) = f^{(j)} \left(\phi^{(t)}, \filter^{(t)} \cdot \pupil^{(p)}, z\right), \\
    & \pupil^{(p)}(\fx, \fy) = \begin{cases}
        1, & \text{if} \, (\fx - \cx{p})^2 + (\fy - \cy{p})^2 \leq r_p^2, \\
        0, & \text{otherwise}
    \end{cases}
\end{split}
\label{eq:pupil}
\end{equation}
%
where $\pupil^{(p)}$ is a binary pupil mask in the Fourier plane, $r_p$ is the radius of the pupil and $\cx{p}, \cy{p}$ are the spatial coordinates of the center of the pupil. This procedure is similar to the pupil-supervision techniques described in \cite{shi2024ergo, schiffers2023stochastic, chakravarthula2022pupil}. Please refer to the supplemental material for more details on our stochastic light field optimization procedure.


\subsection{Implementation Details}

Since we are using a highly-quantized 4-bit phase SLM, the quantization of pixel values needs to be taken into consideration. Such quantization constraints can be enforced using techniques described in prior work~\cite{choi2022flex}. Higher diffraction orders (HDOs) are modeled using the wave propagation model described in \cite{gopakumar2021unfiltered}. We use PyTorch to implement all our algorithms and run optimization.

In all our experiments, the radius $r_p$ of the pupils are set to be 2 mm, resulting in a 4 mm diameter pupil. 81 pupils are equally spaced in the Fourier plane (eyebox plane), where each pupil corresponds to a single view in a $9 \times 9$ light field. The illumination directions of the multisource laser are set such that they match the diffraction angle of the $\pm 1^{\text{st}}$ higher diffraction orders of the blue wavelength. This allows the blue spectrum copies to be perfectly tiled in the Fourier plane, while removing the gaps between the red and green wavelength spectrum copies. Please see the supplemental materials for detailed discussion on the choosing the appropriate illumination angles.

